\chapter{Golden Matter — Particle Masses and Mixing}
\label{ch:20}
% Lane: C — Physics + Data
% Agent: PhysicsWeaver
% Status: COMPLETE

\section{Introduction}

The Standard Model contains 12 fermions: 6 quarks and 6 leptons, each with three generations. This chapter explores geometric relationships between particle masses and mixing matrices and the golden ratio $\phi$.

\section{The 12 Fermion Masses}

The fermion masses span many orders of magnitude:

\begin{table}[h]
\centering
\begin{tabular}{lccc}
\hline
Particle & Mass (GeV) & Generation & Family \\
\hline
\multicolumn{4}{c}{\textit{Quarks}} \\
up ($u$) & 0.00216 & 1 & up-type \\
down ($d$) & 0.00467 & 1 & down-type \\
charm ($c$) & 1.27 & 2 & up-type \\
strange ($s$) & 0.093 & 2 & down-type \\
top ($t$) & 172.76 & 3 & up-type \\
bottom ($b$) & 4.18 & 3 & down-type \\
\hline
\multicolumn{4}{c}{\textit{Leptons}} \\
electron ($e$) & 0.000511 & 1 & charged \\
electron neutrino ($\nu_e$) & $< 2\times10^{-9}$ & 1 & neutral \\
muon ($\mu$) & 0.10566 & 2 & charged \\
muon neutrino ($\nu_\mu$) & $< 2\times10^{-9}$ & 2 & neutral \\
tau ($\tau$) & 1.77686 & 3 & charged \\
tau neutrino ($\nu_\tau$) & $< 2\times10^{-9}$ & 3 & neutral \\
\hline
\end{tabular}
\caption{Standard Model fermion masses (PDG 2024)}
\end{table}

\section{Koide's Lepton Mass Formula}

In 1981, Yoshio Koide discovered a remarkable relation among the charged lepton masses:

\begin{equation}
\frac{m_e + m_\mu + m_\tau}{(\sqrt{m_e} + \sqrt{m_\mu} + \sqrt{m_\tau})^2} = \frac{2}{3}
\end{equation}

This formula, known as the Koide relation, holds to within experimental precision. The value $2/3$ has geometric interpretations related to the golden ratio.

\section{CKM Matrix}

The Cabibbo-Kobayashi-Maskawa (CKM) matrix describes quark mixing in the weak interaction:

\begin{equation}
V_{CKM} = \begin{pmatrix}
V_{ud} & V_{us} & V_{ub} \\
V_{cd} & V_{cs} & V_{cb} \\
V_{td} & V_{ts} & V_{tb}
\end{pmatrix}
\end{equation}

The Jarlskog invariant $J$ quantifies CP violation in the CKM matrix:

\begin{equation}
J \approx 3\times10^{-5}
\end{equation}

\section{PMNS Matrix}

The Pontecorvo-Maki-Nakagawa-Sakata (PMNS) matrix describes neutrino mixing:

\begin{equation}
U_{PMNS} = \begin{pmatrix}
U_{e1} & U_{e2} & U_{e3} \\
U_{\mu1} & U_{\mu2} & U_{\mu3} \\
U_{\tau1} & U_{\tau2} & U_{\tau3}
\end{pmatrix}
\end{pmatrix}

The mixing angles are approximately:
\begin{align}
\theta_{12} &\approx 33.5^\circ \\
\theta_{23} &\approx 49^\circ \\
\theta_{13} &\approx 8.5^\circ
\end{align}

\section{CP Phase $\delta_{CP}$}

The CP-violating phase in the PMNS matrix has been measured to be:

\begin{equation}
\delta_{CP} \approx 195^\circ \approx 1.08\pi
\end{equation}

This value shows deviation from maximal CP violation ($\delta_{CP} = 270^\circ$).

\section{Neutrino Oscillations}

Neutrino oscillations, discovered by Super-Kamiokande (Kajita) and SNO (McDonald) in 1998, earned the 2015 Nobel Prize in Physics. The oscillation probabilities depend on the mass-squared differences:

\begin{align}
\Delta m_{21}^2 &\approx 7.5\times10^{-5} \text{ eV}^2 \\
\Delta m_{31}^2 &\approx 2.5\times10^{-3} \text{ eV}^2
\end{align}

\section{JUNO 2026--2027 Falsification Test}

The Jiangmen Underground Neutrino Observatory (JUNO) in China will measure the neutrino mass ordering with high precision starting in 2026-2027. This experiment will provide crucial tests of neutrino mass models and potential geometric relationships.

\section{Geometric Mass Relations}

The top quark mass shows interesting geometric relationships:

\begin{equation}
m_t \approx 173 \text{ GeV} \approx 100\phi^3 \text{ GeV}
\end{equation}

The weak scale $v \approx 246$ GeV also has geometric interpretations:

\begin{equation}
v \approx \sqrt{2\phi^5} \times 100 \text{ GeV}
\end{equation}

\section{Conclusion}

The pattern of fermion masses and mixing matrices suggests underlying geometric principles. The Koide relation for charged leptons, the structure of mixing matrices, and potential $\phi$-based mass formulas point toward a deeper geometric foundation for particle physics.

% References
% - Koide 1981
% - Kajita-McDonald Nobel 2015
% - JUNO 2026-2027
% - PDG 2024

\begin{thebibliography}{9}
\bibitem{Koide1981} Y. Koide, \textit{A New Formula for the Lepton Mass Sum}, Lett. Nuovo Cimento 34, 1981.
\bibitem{PDG2024} Particle Data Group, \textit{Review of Particle Physics}, Physics Review D 109, 2024.
\end{thebibliography}
